\begin{remark}
\label{remark-cohen-bound-cardinality}
Let $R$ be a ring. Let $\kappa$ be an infinite cardinal.
By applying
Example \ref{example-oka-family-bound-cardinality} and
Proposition \ref{proposition-oka}
we see that any ideal maximal with respect to the property of not being
generated by $\kappa$ elements is prime. This result is not so
useful because there exists a ring for which every prime ideal
of $R$ can be generated by $\aleph_0$ elements, but some
ideal cannot. Namely, let $k$ be a field, let $T$ be a set whose
cardinality is greater than $\aleph_0$ and let
$$
R = k[\{x_n\}_{n \geq 1}, \{z_{t, n}\}_{t \in T, n \geq 0}]/
(x_n^2, z_{t, n}^2, x_n z_{t, n} - z_{t, n - 1})
$$
This is a local ring with unique prime ideal
$\mathfrak m = (x_n)$. But the ideal $(z_{t, n})$ cannot
be generated by countably many elements.
\end{remark}